% Synthetic fixture: named styles declared by \tikzset reach the engine's
% own picture renderer, wherever they were written, and a style applied to
% the picture itself reaches the picture's contents. `ball` is declared in
% the preamble, `slab` beside the picture inside the document body — the
% position a collector that stopped at \begin{document} could not see — and
% both are applied by name on a node. `wire` is applied on the second
% picture's own bracket, so both its edges inherit it, and the second edge
% sets one of the same keys again: the page must carry the edge's own value
% there, the picture's on the rest. The boundary is left open on purpose:
% the census asserts the outlines ship as page paths and that no image box
% stands where the picture is, so the page is the engine's own ink and no
% external tool was asked to draw it. Invented content and design.
\documentclass{article}

\fonts{ dir = "fonts", body = "Source Serif Pro" }

\tikzset{ball/.style={circle, draw, minimum size=7mm}}
\tikzset{wire/.style={draw=blue, thick}}

\begin{document}

\section{Styles the engine reads}

Before the diagram stands a paragraph, so the picture spaces below text.

\tikzset{slab/.style={rectangle, draw, thick, minimum width=9mm, minimum height=6mm}}

\begin{tikzpicture}
  \node[ball] (p) at (0, 0) {P};
  \node[ball] (q) at (3, 0) {Q};
  \node[slab] (r) at (1.5, -1.5) {out};
  \draw (p) -- (q);
\end{tikzpicture}

The next diagram applies its style to itself, not to either edge.

\begin{tikzpicture}[wire]
  \draw (0, 0) -- (3, 0);
  \draw[draw=red] (0, 0.6) -- (3, 0.6);
\end{tikzpicture}

After the diagram stands another, so the flow resumes.

\end{document}
